\documentclass[10pt,letterpaper]{article}
\usepackage{cornell-notes}

\title{Clause 31.11.03.02: Construction And Destruction}
\author{}
\date{}
\setDocTitle{Clause 31.11.03.02: Construction And Destruction}
\setDocSubtitle{C++ 2024}
\setDocOwner{C++ 2024 Cornell Notes}
\setDocDate{\today}
\setCornellCollection{C++ 2024 Cornell Notes}
\setCornellUnitType{Clause}
\setCornellUnitNumber{31.11.03.02}
\setCornellUnitTitle{Construction And Destruction}
\hypersetup{pdftitle={Clause 31.11.03.02: Construction And Destruction},pdfsubject={Cornell notes},pdfauthor={}}

\begin{document}
\maketitle
\begin{CornellOverviewBox}{Overview}
\textbf{Classification:} Standard library - Normative\\[2pt]
\textbf{Learning objective:} Understand the rules, terminology, and semantic consequences associated with Construction and destruction.
\end{CornellOverviewBox}\begin{CornellSummaryBox}{Summary}
{\sffamily\bfseries\color{Accent} Summary}\par\vspace{3pt}Section 31.11.3.2 focuses on Construction and destruction. Apply the specified interface together with its mandates, effects, result conditions, exception behavior, and complexity constraints. When applying it, preserve the distinction between mandatory requirements, recommendations, permissions, and behavior deliberately left undefined, unspecified, or implementation-defined.
\end{CornellSummaryBox}

\end{document}
